% GENERATED FILE. Edit canonical lessons, then run npm run generate:books.
% canonical-source-hash: fnv1a64:ef72a71aa0bb0402
% canonical-lessons: MR-C421-bandhkam, MR-C421-majla, MR-C421-jhopadpatti, MR-C421-ves, MR-C421-hadda

\chapter{Building Sites and Boundaries}
\label{ch:mr-building-sites-and-boundaries}
\hlchaptermodality{421}

\begin{chapteropening}
\textbf{By the end of this chapter:} \emph{I can talk about construction, floors, slums, a village gate and a boundary.}

Complete the last lesson of chapter 421: I can talk about construction, floors, slums, a village gate and a boundary.
\end{chapteropening}

\section[\texorpdfstring{\mr{बांधकाम} (bāndhkām)}{बांधकाम (bāndhkām)}]{\mr{बांधकाम} (bāndhkām) — construction}

\label{lesson:MR-C421-bandhkam}



\begin{quote}
Before the new one: say the Marathi for a resident, then the Marathi for the surroundings.
\end{quote}

\subsection*{You'll want to know: \mr{बांधकाम}}
\textbf{\mr{बांधकाम}} — \emph{bāndhkām} — \textquotedblleft{}construction\textquotedblright{}.

\begin{quote}
\mr{आमच्या} \mr{गल्लीत} \mr{नवीन} \mr{बांधकाम} \mr{चालू} \mr{आहे}.
\end{quote}

\begin{quote}
\emph{āmcyā gallīt navīn bāndhkām cālū āhe.}
\end{quote}

\begin{quote}
New construction is going on in our lane.
\end{quote}

Neuter.

\mr{बांधकाम} \mr{चालू} \mr{आहे}, building work is going on; \mr{बांधकाम} \mr{करणे}, to build.

\begin{cousinweb}
\textbf{\mr{बांधकाम}} (\emph{bāndhkām}) is built from the stem of \textbf{\mr{बांधणे}}, \textquotedblleft{}to tie, to build\textquotedblright{}, and \textbf{\mr{काम}}, \textquotedblleft{}work\textquotedblright{}, inherited from Sanskrit \emph{karman} (CDIAL 2892).
\end{cousinweb}

\begin{grammarlens}[title={What you've built}]
One more word for getting around.
\end{grammarlens}

\subsection*{Your turn}
\begin{itemize}
\raggedright
  \item \emph{Say it:} \emph{bāndhkām}
  \item \emph{Say it:} \emph{bāndhkām}, once more
  \item \emph{Say it:} \emph{āmcyā gallīt navīn bāndhkām cālū āhe.} — New construction is going on in our lane.
  \item \emph{Recall:} say the Marathi for a resident, then the Marathi for the surroundings, then say \emph{bāndhkām} again
\end{itemize}

\begin{tcolorbox}[breakable,colback=teal!4,colframe=teal!35!black,title={Before you move on}]
What does \mr{बांधकाम} mean? (\textquotedblleft{}Construction\textquotedblright{}.) Say it once more.
\end{tcolorbox}
\section[\texorpdfstring{\mr{मजला} (majlā)}{मजला (majlā)}]{\mr{मजला} (majlā) — a floor; a storey}

\label{lesson:MR-C421-majla}



\begin{quote}
Before the new one: say the Marathi for the surroundings, then the Marathi for construction.
\end{quote}

\subsection*{You'll want to know: \mr{मजला}}
\textbf{\mr{मजला}} — \emph{majlā} — \textquotedblleft{}a floor; a storey\textquotedblright{}.

\begin{quote}
\mr{आमचं} \mr{घर} \mr{चौथ्या} \mr{मजल्यावर} \mr{आहे}.
\end{quote}

\begin{quote}
\emph{āmcaṁ ghar cauthyā majalyāvar āhe.}
\end{quote}

\begin{quote}
Our flat is on the fourth floor.
\end{quote}

Masculine.

\mr{तळमजला}, the ground floor; \mr{पहिला} \mr{मजला}, the first floor.

\begin{cousinweb}
\textbf{\mr{मजला}} (\emph{majlā}) is a loan, usually traced to Arabic \emph{manzil}, \textquotedblleft{}a stage of a journey, a halting place\textquotedblright{}, through Persian.
\end{cousinweb}

\begin{grammarlens}[title={What you've built}]
One more word for getting around.
\end{grammarlens}

\subsection*{Your turn}
\begin{itemize}
\raggedright
  \item \emph{Say it:} \emph{majlā}
  \item \emph{Say it:} \emph{majlā}, once more
  \item \emph{Say it:} \emph{āmcaṁ ghar cauthyā majalyāvar āhe.} — Our flat is on the fourth floor.
  \item \emph{Recall:} say the Marathi for the surroundings, then the Marathi for construction, then say \emph{majlā} again
\end{itemize}

\begin{tcolorbox}[breakable,colback=teal!4,colframe=teal!35!black,title={Before you move on}]
What does \mr{मजला} mean? (\textquotedblleft{}A floor; a storey\textquotedblright{}.) Say it once more.
\end{tcolorbox}
\section[\texorpdfstring{\mr{झोपडपट्टी} (jhopaḍpaṭṭī)}{झोपडपट्टी (jhopaḍpaṭṭī)}]{\mr{झोपडपट्टी} (jhopaḍpaṭṭī) — a slum}

\label{lesson:MR-C421-jhopadpatti}



\begin{quote}
Before the new one: say the Marathi for construction, then the Marathi for a floor.
\end{quote}

\subsection*{You'll want to know: \mr{झोपडपट्टी}}
\textbf{\mr{झोपडपट्टी}} — \emph{jhopaḍpaṭṭī} — \textquotedblleft{}a slum\textquotedblright{}.

\begin{quote}
\mr{नदीजवळ} \mr{एक} \mr{मोठी} \mr{झोपडपट्टी} \mr{आहे}.
\end{quote}

\begin{quote}
\emph{nadījavaḷ ek moṭhī jhopaḍpaṭṭī āhe.}
\end{quote}

\begin{quote}
There is a big slum near the river.
\end{quote}

Feminine.

\mr{झोपडी}, a hut, and \mr{पट्टी}, a strip: a row of huts.

\begin{cousinweb}
\textbf{\mr{झोपडपट्टी}} (\emph{jhopaḍpaṭṭī}) is built from \textbf{\mr{झोपडी}}, \textquotedblleft{}a hut\textquotedblright{}, inherited from \emph{*jhuppa} (CDIAL 5403), and \textbf{\mr{पट्टी}}, \textquotedblleft{}a strip\textquotedblright{}.
\end{cousinweb}

\begin{grammarlens}[title={What you've built}]
One more word for getting around.
\end{grammarlens}

\subsection*{Your turn}
\begin{itemize}
\raggedright
  \item \emph{Say it:} \emph{jhopaḍpaṭṭī}
  \item \emph{Say it:} \emph{jhopaḍpaṭṭī}, once more
  \item \emph{Say it:} \emph{nadījavaḷ ek moṭhī jhopaḍpaṭṭī āhe.} — There is a big slum near the river.
  \item \emph{Recall:} say the Marathi for construction, then the Marathi for a floor, then say \emph{jhopaḍpaṭṭī} again
\end{itemize}

\begin{tcolorbox}[breakable,colback=teal!4,colframe=teal!35!black,title={Before you move on}]
What does \mr{झोपडपट्टी} mean? (\textquotedblleft{}A slum\textquotedblright{}.) Say it once more.
\end{tcolorbox}
\section[\texorpdfstring{\mr{वेस} (ves)}{वेस (ves)}]{\mr{वेस} (ves) — a village gate}

\label{lesson:MR-C421-ves}



\begin{quote}
Before the new one: say the Marathi for a floor, then the Marathi for a slum.
\end{quote}

\subsection*{You'll want to know: \mr{वेस}}
\textbf{\mr{वेस}} — \emph{ves} — \textquotedblleft{}a village gate\textquotedblright{}.

\begin{quote}
\mr{गावाच्या} \mr{वेशीवर} \mr{एक} \mr{जुनं} \mr{झाड} \mr{आहे}.
\end{quote}

\begin{quote}
\emph{gāvācyā veśīvar ek junaṁ jhāḍ āhe.}
\end{quote}

\begin{quote}
There is an old tree at the village gate.
\end{quote}

Feminine.

Before an ending \mr{वेस} becomes \mr{वेशी}-: \mr{वेशीवर}, at the gate; also the edge of the village.

\begin{cousinweb}
\textbf{\mr{वेस}} (\emph{ves}) is a Marathi word for the gateway of an old village, and for the village edge; its origin is not settled.
\end{cousinweb}

\begin{grammarlens}[title={What you've built}]
One more word for getting around.
\end{grammarlens}

\subsection*{Your turn}
\begin{itemize}
\raggedright
  \item \emph{Say it:} \emph{ves}
  \item \emph{Say it:} \emph{ves}, once more
  \item \emph{Say it:} \emph{gāvācyā veśīvar ek junaṁ jhāḍ āhe.} — There is an old tree at the village gate.
  \item \emph{Recall:} say the Marathi for a floor, then the Marathi for a slum, then say \emph{ves} again
\end{itemize}

\begin{tcolorbox}[breakable,colback=teal!4,colframe=teal!35!black,title={Before you move on}]
What does \mr{वेस} mean? (\textquotedblleft{}A village gate\textquotedblright{}.) Say it once more.
\end{tcolorbox}
\section[\texorpdfstring{\mr{हद्द} (hadda)}{हद्द (hadda)}]{\mr{हद्द} (hadda) — a boundary; a limit}

\label{lesson:MR-C421-hadda}



\begin{quote}
Before the new one: say the Marathi for a slum, then the Marathi for a village gate.
\end{quote}

\subsection*{You'll want to know: \mr{हद्द}}
\textbf{\mr{हद्द}} — \emph{hadda} — \textquotedblleft{}a boundary; a limit\textquotedblright{}.

\begin{quote}
\mr{नदी} \mr{ही} \mr{दोन} \mr{जिल्ह्यांची} \mr{हद्द} \mr{आहे}.
\end{quote}

\begin{quote}
\emph{nadī hī don jilhyāñcī hadda āhe.}
\end{quote}

\begin{quote}
The river is the boundary between two districts.
\end{quote}

Feminine.

Also a limit of patience: \mr{हद्द} \mr{झाली}, that is the limit.

\begin{cousinweb}
\textbf{\mr{हद्द}} (\emph{hadda}) is a loan through Persian, usually traced to Arabic \emph{ḥadd}, \textquotedblleft{}a limit, a border\textquotedblright{}.
\end{cousinweb}

\begin{grammarlens}[title={What you've built}]
One more word for getting around.
\end{grammarlens}

\subsection*{Your turn}
\begin{itemize}
\raggedright
  \item \emph{Say it:} \emph{hadda}
  \item \emph{Say it:} \emph{hadda}, once more
  \item \emph{Say it:} \emph{nadī hī don jilhyāñcī hadda āhe.} — The river is the boundary between two districts.
  \item \emph{Recall:} say the Marathi for a slum, then the Marathi for a village gate, then say \emph{hadda} again
\end{itemize}

\begin{tcolorbox}[breakable,colback=teal!4,colframe=teal!35!black,title={Before you move on}]
What does \mr{हद्द} mean? (\textquotedblleft{}A boundary; a limit\textquotedblright{}.) Say it once more.
\end{tcolorbox}
